\documentclass{article}
\usepackage{amsmath,amssymb}
\usepackage{xurl}
\usepackage[hidelinks]{hyperref}
% Shared commands for notes in docs/paper-gaps/.

\DeclareUnicodeCharacter{2080}{\ifmmode{}_{0}\else\textsubscript{0}\fi}
\DeclareUnicodeCharacter{2081}{\ifmmode{}_{1}\else\textsubscript{1}\fi}
\DeclareUnicodeCharacter{2082}{\ifmmode{}_{2}\else\textsubscript{2}\fi}
\DeclareUnicodeCharacter{2099}{\ifmmode{}_{n}\else\textsubscript{n}\fi}

\newcommand{\Lean}{\textsc{Lean}}
\ExplSyntaxOn
\NewDocumentCommand{\leanid}{m}{
  \ifmmode
    \textsf{\detokenize{#1}}
  \else
    \group_begin:
    \urlstyle{sf}
    \tl_set:Nn \l_tmpa_tl {#1}
    \tl_replace_all:Nnn \l_tmpa_tl {\_} {_}
    \exp_args:NV \path \l_tmpa_tl
    \group_end:
  \fi
}
\ExplSyntaxOff
\providecommand{\tr}{\operatorname{tr}}
\newcommand{\tnleanIssueBase}{https://github.com/LionSR/TNLean/issues/}
\newcommand{\tnleanPrBase}{https://github.com/LionSR/TNLean/pull/}
\newcommand{\tnleanDocBase}{https://sirui-lu.com/TNLean/docs/}
\newcommand{\ghissue}[1]{\href{\tnleanIssueBase#1}{GitHub issue \##1}}
\newcommand{\ghpr}[1]{\href{\tnleanPrBase#1}{PR \##1}}
\newcommand{\doclink}[1]{\href{\tnleanDocBase#1.html}{\path{#1}}}

% Dirac notation and the identity operator, as in blueprint/src/macros/common.tex.
% \providecommand keeps note-local definitions authoritative.
\usepackage{dsfont}
\providecommand{\ket}[1]{|#1\rangle}
\providecommand{\bra}[1]{\langle #1|}
\providecommand{\braket}[2]{\langle #1 | #2 \rangle}
\providecommand{\ketbra}[2]{|#1\rangle\!\langle #2|}
\providecommand{\Id}{\mathds{1}}
\providecommand{\one}{\mathds{1}}
\providecommand{\C}{\mathbb{C}}
\providecommand{\MN}[1]{M_{#1}(\C)}
\providecommand{\MD}{M_D(\C)}

% Verdict marker: \gapnote{<kind>}{<status>}. Kind is the policy
% classification (clarification, local-correction, scope-restriction,
% unfaithful, false-source, open-gap); status is open, wip (elimination
% underway), resolved, or historical. Machine-read by the site index;
% prints a verdict line.
\newcommand{\gapnote}[2]{\par\noindent\textbf{Verdict:} #1 (#2).\par}

\title{Square supplied gates in the printed MPU blocking argument}
\author{}
\date{2026-11-16}
\begin{document}
\maketitle
\gapnote{scope-restriction}{open}

\section{Source and the proved construction}

The proof of Proposition~\texttt{prop:sigma1-and-sigma2-equal} in
\cite[lines~1643--1733]{Cirac2017MPU} chooses particular blocked gates,
\[
    u^{(k)}=(I\otimes v^{\otimes(k-1)}\otimes I)u^{\otimes k},\qquad
    v^{(k)}=I^{\otimes(k-1)}\otimes v\otimes I^{\otimes(k-1)}.
\]
The construction
\leanid{MPOTensor.TwoSiteStandardFormData.printedBlock} proves that these
are the gates of the actual positive-length blocked tensor, after its
explicit physical half-word relabelling. It retains the original bond
space and constructs both half factors by adjoining identity spectators.
Its input is supplied square standard-form data: both gates act on two
copies of the original spin space. No canonical ambient presentation,
comparison gauge, symmetry, or previously supplied blocked factors are
premises of this construction.

The theorem
\leanid{MPOTensor.TwoSiteStandardFormData.printedBlocking_trace_eq_intrinsic}
identifies the normalized trace of these particular gates with the intrinsic
trace of the actual blocked tensor. It uses the constructed open standard
form and invariance of the two-ring exchange trace under a common physical
basis relabelling.
The theorem
\leanid{MPOTensor.TwoSiteStandardFormData.printed_blocking_trace_parity}
then proves the odd/even trace identities under adjoint symmetry on every
even ring of length greater than one. It consumes the existing intrinsic
parity theorem rather than assuming the blocked traces or their parity.

\section{Remaining source integration}

A general supplied standard form permits distinct intermediate dimensions
$\ell,r$ and rectangular gates between the physical and intermediate pair
spaces. The new construction starts in the square case $\ell=r=d$.
The source time-reversal argument first obtains this situation from its
index and local-symmetry results. The new theorem does not derive that
reduction from an arbitrary supplied rectangular presentation, nor choose
a source standard form for an arbitrary MPU. Those are separate results
with existing owners. The original unrestricted source-labelled parity
entries therefore remain unfinished.

The local-unitary and sequential-unwinding corollaries that start from
\leanid{MPOTensor.IsMPUCanonicalFormII} retain their explicitly documented
canonical ambient restriction; see the full-support note. That hypothesis
is not present in the new actual blocked supplier or its trace parity.
The latter asserts identities of normalized traces. It does not by itself
supply every local time-reversal certificate or replace the general
local-symmetry criterion.

Removing this scope distinction requires integrating the general
standard-form and time-reversal rank reductions, transporting the actual
printed gates through their specified physical identifications, and
checking the resulting source statement and all review gates. No
continuous gauge choice is part of this blocking construction.

\bibliographystyle{alpha}
\bibliography{references}
\end{document}
